% The .s1 format (main text). The normative specification is a separate supplementary file (spec.tex).
\section{The \sone{} format}
\label{sec:format}\label{sec:how}

Fig.~\ref{fig:how} shows the format from a whole dataset down to individual bits; the full specification is a separate
supplementary file. A dataset is a small manifest that lists one file per sample, each named by a cryptographic hash of
its contents, so a link always refers to exactly the same bytes and anyone can check them (Fig.~\ref{fig:how}a).
Information shared by many samples, such as the gene names of a genome annotation, is stored once and downloaded once
(Fig.~\ref{fig:how}b). Table~\ref{tab:props} compares these properties with those of other formats.

\para{Blocks for partial reads}How does a reader get a few hundred cells out of a remote file without reading the rest?
Each count matrix starts with a small index that says where every block of consecutive cells lives in the file, and each
block decodes on its own (Fig.~\ref{fig:how}c). A reader fetches the index once and then only the blocks it needs, with one
request per run of blocks, from a local disk, a web server or a cloud bucket alike. Because every block is located before
any is read, all the requests of a minibatch can be sent at once.

\para{A model of the counts}Inside a block, genes are numbered from most to least frequently detected, so each cell's
non-zero entries become a sequence of small steps and small counts, most of them ones. Each entry becomes one short
symbol, and its probability is looked up in a table chosen by two things the reader already knows: how many entries the
cell has, and how far along the gene list it is (Fig.~\ref{fig:how}d). An rANS entropy coder~\cite{duda2013asymmetric}
then spends close to the minimum number of bits on each symbol. The tables are fixed for the whole matrix, so decoding is
a few table look-ups per entry, and 32 coders run side by side so that a processor with AVX2 decodes eight to sixteen
symbols at once. Writing a file ends with a full decode that must reproduce the input exactly.

% Figure 2: how .s1 stores single-cell counts, as four panels joined by zoom wedges, lettered as sections/format.tex
% cites them: (a) a dataset, (b) one sample file (header, member table, shared packs), (c) one count-matrix member
% (prefix, chunks, range requests), (d) one chunk. Plain words stand in for spec terms: "table of contents" is the
% member table, "block" a chunk, "code tables" the entropy tables, "gene order" the row keys, "symbol" the gap symbol
% (sigma, amber: the residual rho), "coders" the rANS lanes. Sizes, the prefix, the request pattern, the 12 x 40
% contexts (ctx_kind 0) and the chunk byte order follow sections/supplement_spec.tex. Genes, counts and ranks in (d)
% are a worked example of those rules, not data. Request labels in (c) follow the reference client
% (python/singlet/s1/client.py: the first request carries the prefix, then one request per run of chunks).
\begin{figure*}[tbp]
\centering
% panel title with a hanging indent after the panel letter
\def\ptitle#1#2#3#4{\node[ptit, text width=#3] at (#1) {\hyphenpenalty=10000\exhyphenpenalty=10000\hangindent=0.36cm\hangafter=1\makebox[0.36cm][l]{\small #2}#4};}
% station title in panel d
\def\stitle#1#2#3#4{\node[stit, text width=#2] at (#1) {\hyphenpenalty=10000\exhyphenpenalty=10000\hangindent=0.28cm\hangafter=1\makebox[0.28cm][l]{\bfseries #3}#4};}
% zoom wedge from a parent element's corners (#1, #2) to a child panel's corners (#3, #4)
\def\zoomwedge#1#2#3#4{\begin{scope}[on background layer]\fill[sone!5] (#1)--(#3)--(#4)--(#2)--cycle;\draw[wl] (#1)--(#3) (#2)--(#4);\end{scope}}
% one row of the table of contents in (b): y, style, name, start, length
\def\tocrow#1#2#3#4#5{%
  \node[cell,#2,minimum width=1.95cm,anchor=north west] at (6.75,#1) {#3};
  \node[cell,#2,minimum width=0.80cm,anchor=north west,align=center] at (8.70,#1) {#4};
  \node[cell,#2,minimum width=0.80cm,anchor=north west,align=center] at (9.50,#1) {#5};}
% legend chips for the note in (d): blue symbol, amber symbol
\def\chipkey#1#2{{\setlength{\fboxsep}{0pt}\fcolorbox{#1}{#2}{\rule{0pt}{5.5pt}\hspace{8pt}}}}
\begin{tikzpicture}[
  x=1cm, y=1cm,
  font=\sffamily\scriptsize,
  >={Stealth[length=3pt,width=2.4pt]},
  st/.style={draw=sone!80!black, fill=sone!14, line width=0.4pt},
  hot/.style={draw=sone!80!black, fill=sone!75, text=white, line width=0.4pt},
  idle/.style={draw=gray!55, fill=gray!15, line width=0.4pt},
  ghost/.style={draw=gray!60, dashed, fill=white, text=ink2, line width=0.4pt},
  amb/.style={draw=dec!80!black, fill=dec!30, line width=0.4pt},
  seg/.style={rectangle, inner sep=3pt, outer sep=0pt, align=center, anchor=south west},
  box/.style={rectangle, inner sep=4pt, outer sep=0pt, align=center, rounded corners=1pt},
  cell/.style={rectangle, inner sep=3pt, outer sep=0pt, align=left, minimum height=0.36cm, text depth=0.5pt},
  frame/.style={draw=gray!35, rounded corners=2pt, line width=0.5pt},
  ptit/.style={font=\sffamily\scriptsize\bfseries, anchor=north west, inner sep=0pt, align=left},
  stit/.style={font=\sffamily\scriptsize, anchor=north west, inner sep=0pt, align=left},
  note/.style={text=ink2, inner sep=0pt, align=left},
  hd/.style={font=\sffamily\tiny, text=ink2, inner sep=0pt},
  wl/.style={draw=gray!55, line width=0.4pt},
  arr/.style={->, ink2, line width=0.45pt},
  chip/.style={rectangle, inner sep=0pt, outer sep=0pt, minimum height=0.27cm, minimum width=0.36cm, line width=0.4pt},
]

% =============================================================== (a) a dataset
\ptitle{0,0}{a}{4.1cm}{A dataset}
\node[box, st, anchor=north west, minimum width=1.75cm, minimum height=1.84cm] (man) at (0,-0.62)
  {\textbf{manifest}\\[2pt]lists every\\sample file};
\node[box, hot, anchor=north west, minimum width=1.6cm, minimum height=0.36cm] (s1) at (2.2,-0.62) {sample 1\textbf{.s1}};
\node[box, st,  anchor=north west, minimum width=1.6cm, minimum height=0.36cm] (s2) at (2.2,-1.18) {sample 2\textbf{.s1}};
\node at (3.0,-1.83) {$\vdots$};
\node[box, st,  anchor=north west, minimum width=1.6cm, minimum height=0.36cm] (sN) at (2.2,-2.04) {sample \textit{N}\textbf{.s1}};
\foreach \s in {s1,s2,sN} \draw[arr] (man.east) -- (\s.west);
\node[note, anchor=north west, text width=3.75cm] at (0,-2.72)
  {Each file is named by a hash of its bytes, so a link always means the same data.};

% =============================================================== (b) one sample file
\draw[frame] (4.5,0.02) rectangle (16.4,-3.4);
\ptitle{4.62,-0.08}{b}{6cm}{One sample file}
% the file as stored
\node[seg, st,  minimum width=0.95cm, minimum height=0.62cm] (hdrB) at (4.65,-3.22) {header};
\node[seg, st,  minimum width=1.30cm, minimum height=0.62cm] (tocB) at (hdrB.south east) {table of\\contents};
\node[seg,      minimum width=0.40cm, minimum height=0.62cm, text=gray] (brkB) at (tocB.south east) {//};
\node[seg, st,  minimum width=1.55cm, minimum height=0.62cm] (celB) at (brkB.south east) {cells\\table};
\node[seg, hot, minimum width=1.80cm, minimum height=0.62cm] (cntB) at (celB.south east) {\textbf{exon counts}};
\node[seg, st,  minimum width=1.80cm, minimum height=0.62cm] (intB) at (cntB.south east) {intron counts};
\node[seg, ghost, minimum width=2.15cm, minimum height=0.62cm] (derB) at (intB.south east) {cell calls\\recipe, 0 bytes};
% zoom: header and table of contents
\node[box, st, anchor=north west, minimum width=1.85cm, minimum height=1.66cm] (hdrZ) at (4.65,-0.5)
  {\textbf{header}\\[1pt]{\tiny 64 bytes}\\[2pt]file type\\sample ID\\file length};
\node[cell, anchor=north west, minimum width=1.95cm, font=\sffamily\scriptsize\bfseries] at (6.75,-0.5) {table of contents};
\node[hd, anchor=north] at (9.10,-0.6) {start};
\node[hd, anchor=north] at (9.90,-0.6) {length};
\tocrow{-0.83}{st}{cells}{\textit{o}\textsubscript{1}}{\textit{l}\textsubscript{1}}
\tocrow{-1.19}{hot}{exon counts}{\textit{o}\textsubscript{2}}{\textit{l}\textsubscript{2}}
\tocrow{-1.55}{st}{intron counts}{\textit{o}\textsubscript{3}}{\textit{l}\textsubscript{3}}
\tocrow{-1.91}{ghost}{cell calls}{0}{0}
\zoomwedge{hdrB.north west}{tocB.north east}{4.65,-2.24}{10.3,-2.24}
% shared packs
\node[note, anchor=north west, text width=3.8cm, text=black] at (12.3,-0.5)
  {\textbf{shared packs}\\\textcolor{ink2}{stored once, downloaded once}};
\node[box, ghost, anchor=north west, minimum width=1.8cm, minimum height=0.5cm, text height=1.6ex, text depth=0.45ex] (pkG) at (12.3,-1.3) {gene names};
\node[box, ghost, anchor=north west, minimum width=1.9cm, minimum height=0.5cm, text height=1.6ex, text depth=0.45ex] (pkB) at (14.3,-1.3) {barcode lists};
\draw[->, dashed, ink2, line width=0.4pt] ($(cntB.north east)+(-0.35,0)$) -- ($(pkG.south west)+(0.35,0)$);
\draw[->, dashed, ink2, line width=0.4pt] ($(intB.north)+(0.3,0)$) -- ($(pkG.south)+(0.1,0)$);

% =============================================================== (c) one count matrix
\draw[frame] (0,-3.9) rectangle (16.4,-6.7);
\ptitle{0.12,-4.0}{c}{12cm}{One count matrix: a small prefix, then blocks of cells}
\node[seg, hot, minimum width=0.95cm, minimum height=0.72cm] (hdrC) at (0.15,-6.1) {header};
\node[seg, hot, minimum width=2.30cm, minimum height=0.72cm] (idxC) at (hdrC.south east) {};
\node[text=white, inner sep=0pt] at ($(idxC.north)+(0,-0.2)$) {block index};
\foreach \i in {0,...,8} \draw[white, line width=0.3pt] ($(idxC.south west)+(0.25+0.21*\i,0.1)$) rectangle ++(0.15,0.15);
\node[seg, hot, minimum width=1.30cm, minimum height=0.72cm] (tabC) at (idxC.south east) {code\\tables};
\node[seg, hot, minimum width=1.20cm, minimum height=0.72cm] (ordC) at (tabC.south east) {gene\\order};
\node[seg, st,  minimum width=0.85cm, minimum height=0.72cm] (b1) at (ordC.south east) {block\\1};
\node[seg, hot, minimum width=0.85cm, minimum height=0.72cm] (b2) at (b1.south east) {block\\2};
\node[seg, hot, minimum width=0.85cm, minimum height=0.72cm] (b3) at (b2.south east) {block\\3};
\node[seg, st,  minimum width=0.55cm, minimum height=0.72cm] (bd1) at (b3.south east) {$\cdots$};
\node[seg, hot, minimum width=0.85cm, minimum height=0.72cm] (bk) at (bd1.south east) {block\\\textit{k}};
\node[seg, st,  minimum width=0.55cm, minimum height=0.72cm] (bd2) at (bk.south east) {$\cdots$};
\node[seg, st,  minimum width=0.85cm, minimum height=0.72cm] (bK) at (bd2.south east) {block\\\textit{K}};
\node[seg, idle, minimum width=1.30cm, minimum height=0.72cm] (nmC) at (bK.south east) {names};
% the prefix
\draw[decorate, decoration={brace, amplitude=3pt}, ink2, line width=0.4pt] ($(hdrC.north west)+(0.02,0.06)$) -- ($(ordC.north east)+(-0.02,0.06)$)
  node[midway, above=5pt, note, align=center, text=black] {\textbf{prefix} (request 1): read once,\\enough to plan any read};
% the wanted cells and one request per run of blocks
\foreach \x in {7.20,7.33,7.47,7.62,7.80,7.95,8.10,8.32,9.55,9.80}
  \draw[sone!85!black, line width=0.8pt] (\x,-5.33) -- (\x,-5.19);
\draw[decorate, decoration={brace, amplitude=2.5pt}, ink2, line width=0.4pt] ($(b2.north west)+(0.02,0.3)$) -- ($(b3.north east)+(-0.02,0.3)$)
  node[midway, above=4pt, inner sep=0pt] {request 2};
\draw[decorate, decoration={brace, amplitude=2.5pt}, ink2, line width=0.4pt] ($(bk.north west)+(0.02,0.3)$) -- ($(bk.north east)+(-0.02,0.3)$)
  node[midway, above=4pt, inner sep=0pt] {request 3};
% index entries point at block starts
\draw[->, sone!80!black, line width=0.35pt] (2.2,-6.1) .. controls (3.2,-6.45) and (6.2,-6.45) .. ($(b2.south west)+(0.04,0)$);
\draw[->, sone!80!black, line width=0.35pt] (2.2,-6.1) .. controls (3.4,-6.6) and (7.0,-6.6) .. ($(b3.south west)+(0.04,0)$);
% legend
\draw[sone!85!black, line width=0.9pt] (13.46,-4.64) -- (13.46,-4.46) (13.56,-4.64) -- (13.56,-4.46);
\node[anchor=west, inner sep=0pt] at (13.75,-4.55) {cells you want};
\draw[hot] (13.4,-5.04) rectangle ++(0.22,0.22);
\node[anchor=west, inner sep=0pt] at (13.75,-4.93) {fetched};
\draw[st] (13.4,-5.44) rectangle ++(0.22,0.22);
\node[anchor=west, inner sep=0pt] at (13.75,-5.33) {not fetched};
\node[note, anchor=north west, text width=2.9cm] at (13.4,-5.7) {Blocks hold consecutive cells; each decodes alone.};

% =============================================================== (d) inside one block
\draw[frame] (0,-7.2) rectangle (16.4,-12.2);
\ptitle{0.12,-7.3}{d}{16cm}{Inside one block: each cell becomes small numbers, then short codes}
% --- station 1: one cell's counts
\stitle{0.15,-7.72}{2.3cm}{1}{One cell's counts}
\node[hd, anchor=east] at (1.3,-8.46) {gene};
\node[hd] at (1.75,-8.46) {count};
\foreach \g/\v [count=\i from 0] in {ACTB/1,CD3E/2,MALAT1/3,MS4A1/1,RPL13/1} {
  \node[anchor=east, inner sep=0pt] at (1.3,{-8.74-0.34*\i}) {\g};
  \node[chip, st] at (1.75,{-8.74-0.34*\i}) {\v};}
\node[note, anchor=north west] at (0.15,-10.42) {zeros not stored};
\draw[arr] (2.4,-9.42) -- (2.75,-9.42);
% --- station 2: genes numbered by detection
\stitle{2.9,-7.72}{3.6cm}{2}{Genes numbered, most often detected first}
\node[hd, anchor=east] at (3.95,-8.46) {gene};
\node[hd, anchor=west] at (4.05,-8.46) {cells detecting it};
\node[hd] at (5.85,-8.46) {no.};
\node[hd] at (6.3,-8.46) {count};
\foreach \g/\w/\d/\v [count=\i from 0] in {MALAT1/1.45/0/3,ACTB/1.3/1/1,RPL13/1.15/2/1,CD3E/0.55/9/2,MS4A1/0.25/31/1} {
  \node[anchor=east, inner sep=0pt] at (3.95,{-8.74-0.34*\i}) {\g};
  \fill[gray!45] (4.05,{-8.80-0.34*\i}) rectangle ++(\w,0.12);
  \node[font=\sffamily\scriptsize\bfseries, text=sone!80!black, inner sep=0pt] at (5.85,{-8.74-0.34*\i}) {\d};
  \node[chip, st] at (6.3,{-8.74-0.34*\i}) {\v};}
\node[note, anchor=north west, text width=3.6cm] at (2.9,-10.42) {one numbering for the whole matrix: the gene order in the prefix};
\draw[arr] (6.6,-9.42) -- (6.95,-9.42);
% --- station 3: steps and counts become symbols
\stitle{7.1,-7.72}{3.0cm}{3}{Small steps and small counts}
\node[hd] at (7.3,-8.46) {no.};
\node[hd] at (7.8,-8.46) {step};
\node[hd] at (8.33,-8.46) {count};
\node[hd, anchor=west] at (8.72,-8.46) {symbols};
\foreach \d/\s/\v/\two [count=\i from 0] in {0/0/3/1,1/0/1/0,2/0/1/0,9/6/2/1,31/21/1/0} {
  \node[inner sep=0pt] at (7.3,{-8.74-0.34*\i}) {\d};
  \node[inner sep=0pt] at (7.8,{-8.74-0.34*\i}) {\s};
  \node[inner sep=0pt] at (8.33,{-8.74-0.34*\i}) {\v};
  \node[chip, st, anchor=west] at (8.72,{-8.74-0.34*\i}) {};
  \ifnum\two=1 \node[chip, amb, anchor=west] at (9.14,{-8.74-0.34*\i}) {};\fi}
\node[note, anchor=north west, text width=3.0cm] at (7.1,-10.42)
  {step = genes skipped\\[1pt]
   \chipkey{sone!80!black}{sone!14}\ one per gene\\[1pt]
   \chipkey{dec!80!black}{dec!30}\ one more if count $\neq$\,1};
\draw[arr] (9.7,-9.42) -- (10.05,-9.42);
% --- station 4: a code table chosen by context
\stitle{10.2,-7.72}{3.0cm}{4}{Short codes for common symbols}
\begin{scope}[shift={(11.2,-9.8)}]
  \fill[sone!14] (0,0) rectangle (2.0,1.08);
  \draw[white, line width=0.15pt, xstep=0.05, ystep=0.09] (0,0) grid (2.0,1.08);
  \draw[sone!80!black, line width=0.4pt] (0,0) rectangle (2.0,1.08);
  \draw[sone!80!black, line width=0.9pt] (5*0.05,2*0.09) rectangle ++(0.05,0.09);
\end{scope}
\node[hd, align=right, anchor=east] at (11.12,-9.26) {genes\\detected\\in cell:\\12 classes};
\node[hd, anchor=north, align=center] at (12.2,-9.88) {position in gene list:\\40 buckets};
\node[note, anchor=north west, text width=3.1cm] at (10.2,-10.42)
  {one table per box, stored once in the prefix; common symbols cost few bits};
\draw[arr] (13.25,-9.42) -- (13.6,-9.42);
% --- station 5: 32 coders, one word stream
\stitle{13.75,-7.72}{2.55cm}{5}{32 rANS coders side by side (8 shown)}
\foreach \l in {0,...,7} {
  \draw[sone!70, line width=0.8pt] (14.3,{-8.74-0.13*\l}) -- (14.95,{-8.74-0.13*\l});
  \draw[sone!45, line width=0.4pt] (14.95,{-8.74-0.13*\l}) -- (15.25,-9.195);}
\node[hd, anchor=east] at (14.25,-8.74) {coder 1};
\node[hd, anchor=east] at (14.25,-9.65) {coder 8};
\foreach \w in {0,...,3} \draw[st] (15.3+0.22*\w,-9.295) rectangle ++(0.2,0.2);
\node[hd, anchor=north] at (15.73,-9.42) {16-bit words};
\node[note, anchor=north west, text width=2.55cm] at (13.75,-9.98)
  {AVX2 decodes 8 at once, AVX-512 16. Each coder must end where encoding began: a \mbox{built-in} check.};
% --- the block as stored
\node[anchor=west, inner sep=0pt, text=ink2] at (0.15,-11.83) {stored as};
\node[seg, st, minimum width=0.85cm, minimum height=0.46cm, text height=1.6ex, text depth=0.45ex] (f1) at (1.35,-12.06) {sizes};
\node[seg, st, minimum width=2.9cm, minimum height=0.46cm, text height=1.6ex, text depth=0.45ex] (f2) at (f1.south east) {genes detected per cell};
\node[seg, st, minimum width=2.1cm, minimum height=0.46cm, text height=1.6ex, text depth=0.45ex] (f3) at (f2.south east) {coder start points};
\node[seg, st, minimum width=2.3cm, minimum height=0.46cm, text height=1.6ex, text depth=0.45ex] (f4) at (f3.south east) {32 coder states};
\node[seg, st, minimum width=6.75cm, minimum height=0.46cm, text height=1.6ex, text depth=0.45ex] (f5) at (f4.south east) {one stream of 16-bit words, shared by all coders};

% =============================================================== zoom wedges
\zoomwedge{cntB.south west}{cntB.south east}{0,-3.9}{16.4,-3.9}
\zoomwedge{bk.south west}{bk.south east}{0,-7.2}{16.4,-7.2}
\zoomwedge{s1.north east}{s1.south east}{4.5,0.02}{4.5,-3.4}
\end{tikzpicture}
\caption{\textbf{How \sone{} stores single-cell counts.}
\textbf{(a)} A dataset lists per-sample files, each named by a hash of its bytes.
\textbf{(b)} A header and table of contents locate each part; gene names and barcodes sit in shared packs, and
derivable tables are recipes verified on that sample.
\textbf{(c)} A small prefix plans any read, so a reader fetches only the blocks (chunks) of cells it needs.
\textbf{(d)} Genes numbered by detection turn each cell into small steps and counts, coded in few bits with tables
chosen by context.}
\label{fig:how}
\end{figure*}

% Main-text table: what a stored count matrix offers besides size and read speed.
\begin{table*}[tbp]
\centering
\caption{Properties beyond size and speed, as each format is specified}
\label{tab:props}
\footnotesize
\begin{tabularx}{\linewidth}{@{}>{\raggedright}p{2.7cm} >{\raggedright}X >{\raggedright}X c >{\raggedright}p{3.3cm} c@{}}
\toprule
Format & Reads in part over HTTP & Integrity & Digest & Several matrices + cell table & Integers \tabularnewline
\midrule
\textbf{\sone} & yes: block index in the prefix, one range per run of cells & object named by its SHA-256 prefix; header CRC-32; digest per member & yes & yes, plus recipes and shared packs & yes \tabularnewline
h5ad, 10x HDF5, loom & with chunked datasets and an HDF5 range driver & Fletcher-32 per chunk (opt.) & no & yes (h5ad, loom) & no \tabularnewline
AnnData Zarr & yes: one object per chunk & CRC32C codec (opt., Zarr v3) & no & yes & no \tabularnewline
TileDB-SOMA & yes, through the TileDB library & checksum filters (opt.) & no & yes & no \tabularnewline
Parquet COO & yes: row groups and pages & page CRC-32 (opt.) & no & no & no \tabularnewline
BPCells & no (local directory) & no & no & no & no \tabularnewline
rds, MTX.gz & no (whole file) & gzip/xz stream checksums & no & no (rds: one object) & no \tabularnewline
\bottomrule
\end{tabularx}
\par\vspace{3pt}
\begin{minipage}{\linewidth}\scriptsize
Reads in part over HTTP: some cells by range request from plain storage, with no server software. Integrity: checks a
reader can apply to the bytes. Digest: a hash of the content (values, names) independent of the encoding. Integers: the
format or its writer rejects matrices that are not non-negative integer counts. opt.: an option writers rarely enable.
\end{minipage}
\end{table*}


\para{Stored gene order}A reader receives each cell's entries in the stored, frequency-sorted gene order, with one
permutation of the genes for the whole matrix. Returning them in the source order would mean sorting every cell's entries,
which costs \nm{s1order-x} times the decoding time on one core; a training loader instead applies the permutation once, to
its gene vocabulary, and matrices from many studies are mapped onto a shared gene axis the same way. All results here read
\sone{} in its stored order.
